% !TEX root = ../main.tex

\section{第 15.1 节课程题目：总自旋与超精细分裂}
\label{exercises:QM-C15-S01}
\exercisemeta
  {QM-C15-S01}
  {2026-11-01 至 2026-11-02}
  {Shankar, Second Edition；印刷 pp.~405--407 / PDF pp.~411--413}
  {正文示范已讲解；15.1.1 关键作用已作答、其余补全；15.1.2 第 (1) 问已作答，均已核对}
  {笔记已确认}

本次选做 15.1.1 和 15.1.2 第 (1) 问；15.1.2 第 (2)、(3) 问未选做。

\exercisequestion{QM-C15-S01-E01}{教材正文示范：氦原子基态的自旋结构}
\textbf{来源：}印刷 pp.~406--407，式 (15.1.17)--(15.1.19)，无独立 Example 题号。

\textbf{对应知识点：}\relatedtopic{QM-C15-S01-T03}{交换对称性}。

\textbf{问题。}在独立粒子图像中，氦原子基态的两个电子均占据同一个归一化 $1s$ 空间轨道 $u$。构造满足费米子交换条件的归一化完整态，并确定总轨道角动量与总自旋。

\begin{checkedsolution}
空间部分 $\Phi(\mathbf r_1,\mathbf r_2)=u(\mathbf r_1)u(\mathbf r_2)$ 已归一化且对称，因此自旋部分须取唯一的反对称态：
\begin{equation}
  \boxed{\Psi(\mathbf r_1,\mathbf r_2)
  =u(\mathbf r_1)u(\mathbf r_2)
    \frac{\ket{+-}-\ket{-+}}{\sqrt2}.}
\end{equation}
交换全部标签后 $\Psi\mapsto-\Psi$。两个 $1s$ 因子分别有 $l=0$，自旋部分有 $s=0$，故 $L^2\Psi=S^2\Psi=0$。自旋态是单重态，不是单独一个 $\ket{+-}$；空间乘积采用独立粒子近似，包含电子关联时仍须满足交换条件。
\end{checkedsolution}

\exercisequestion{QM-C15-S01-E02}{教材习题 15.1.1：构造总自旋平方矩阵}
\textbf{来源与知识点：}印刷 p.~405；\relatedtopic{QM-C15-S01-T01}{乘积基}与\relatedtopic{QM-C15-S01-T02}{总自旋基}。

\textbf{题目。}两个粒子均有自旋 $1/2$。按 $\mathcal B_p=(\ket{++},\ket{+-},\ket{-+},\ket{--})$ 的顺序，推导教材式 (15.1.10) 的 $S^2$ 矩阵，并求出 $S_z=0$ 子空间中使 $S^2$ 对角化的两个归一化组合，即式 (15.1.11)。可用
\begin{equation}
  \mathbf S_1\cdot\mathbf S_2
  =S_{1z}S_{2z}+\frac12(S_{1+}S_{2-}+S_{1-}S_{2+}).
\end{equation}

\begin{checkedsolution}
\textbf{1. 最关键的一列。}展开
\begin{equation}
  S^2=S_1^2+S_2^2+2S_{1z}S_{2z}
     +S_{1+}S_{2-}+S_{1-}S_{2+}.
\end{equation}
作用于 $\ket{+-}$，各项分别为
\begin{align}
  (S_1^2+S_2^2)\ket{+-}&=\frac32\hbar^2\ket{+-},\\
  2S_{1z}S_{2z}\ket{+-}&=-\frac12\hbar^2\ket{+-},\\
  S_{1+}S_{2-}\ket{+-}&=0,\qquad
  S_{1-}S_{2+}\ket{+-}=\hbar^2\ket{-+}.
\end{align}
所以
\begin{equation}
  \boxed{S^2\ket{+-}=\hbar^2(\ket{+-}+\ket{-+}).}
\end{equation}
非对角项来自升降算符的交叉作用，因而 $\ket{+-}$ 自身不是 $S^2$ 本征态。

\textbf{2. 其余列。}交换标签得到 $S^2\ket{-+}=\hbar^2(\ket{+-}+\ket{-+})$。
对两个端点态，升降交叉项为零，$2S_{1z}S_{2z}$ 均贡献 $\hbar^2/2$，所以 $S^2\ket{++}=2\hbar^2\ket{++}$、$S^2\ket{--}=2\hbar^2\ket{--}$。把作用结果依次作为列：
\begin{equation}
  \boxed{S^2\ \longleftrightarrow\ \hbar^2
  \begin{pmatrix}
    2&0&0&0\\
    0&1&1&0\\
    0&1&1&0\\
    0&0&0&2
  \end{pmatrix}.}
\end{equation}

\textbf{3. 对角化中间子空间。}在 $(\ket{+-},\ket{-+})$ 中，矩阵是 $\hbar^2\begin{psmallmatrix}1&1\\1&1\end{psmallmatrix}$。无量纲本征值满足 $(1-\lambda)^2-1=\lambda(\lambda-2)=0$。
$\lambda=2$ 要求两个系数相等，$\lambda=0$ 要求系数相反。归一化后
\begin{equation}
  \boxed{\ket{1,0}=\frac{\ket{+-}+\ket{-+}}{\sqrt2},\qquad
    \ket{0,0}=\frac{\ket{+-}-\ket{-+}}{\sqrt2}.}
\end{equation}
它们分别满足 $S^2=2\hbar^2$、$S^2=0$，总投影都为零；加上两个端点态，四个正交态完整覆盖原空间。
\end{checkedsolution}

\exercisequestion{QM-C15-S01-E03}{教材习题 15.1.2 第 (1) 问：氢原子超精细分裂}
\textbf{来源与知识点：}印刷 p.~407；\relatedtopic{QM-C15-S01-T04}{自旋耦合与基的选择}。

\textbf{题目。}将氢原子空间轨道固定为 $1s$，未扰动能量为 $-\mathrm{Ry}$。电子和质子均有自旋 $1/2$，加入
\begin{equation}
  H_{\rm hf}=A\mathbf S_e\cdot\mathbf S_p,\qquad A>0,
  \qquad H=-\mathrm{Ry}\,I+H_{\rm hf}.
\end{equation}
求两组能量及对应的总自旋态，确定哪组较低、各自简并度和能级间距。不加外磁场。

\begin{checkedsolution}
由 $\mathbf S_e\cdot\mathbf S_p=(S^2-S_e^2-S_p^2)/2$ 和 $S_e^2=S_p^2=3\hbar^2I/4$，在 $\ket{s,m}$ 上取本征值：
\begin{equation}
  E(s)=-\mathrm{Ry}+\frac{A\hbar^2}{2}
    \left[s(s+1)-\frac32\right].
\end{equation}
\begin{center}
\begin{tabular}{@{}llll@{}}
  \toprule
  类型 & $s$ & 能量 & 简并度 \\
  \midrule
  三重态 & $1$ & $E_{\rm t}=-\mathrm{Ry}+A\hbar^2/4$ & $3$ \\
  单重态 & $0$ & $E_{\rm s}=-\mathrm{Ry}-3A\hbar^2/4$ & $1$ \\
  \bottomrule
\end{tabular}
\end{center}
$A>0$，故单重态较低，能级间距为
\begin{equation}
  \boxed{\Delta E=E_{\rm t}-E_{\rm s}=A\hbar^2.}
\end{equation}
三重态的三个 $m$ 同能，因为此模型只依赖 $S^2$。共同的库仑能量在能级差中抵消；$A\hbar^2$ 的量纲为能量。
\end{checkedsolution}

\section{第 15.2 节课程题目：CG 展开与子空间投影}
\label{exercises:QM-C15-S02}
\exercisemeta
  {QM-C15-S02}
  {2026-11-02 至 2026-11-04}
  {Shankar, Second Edition；正文换基示范在印刷 pp.~413--414 / PDF pp.~419--420；习题在印刷 pp.~413、415 / PDF pp.~419、421}
  {15.2.2 第 (1) 问的两个关键态、15.2.5 已作答并核对；其余 CG 态已补全}
  {笔记已确认}

本次选做 15.2.2 第 (1) 问和 15.2.5。另收录教材正文的换基矩阵示范，单重态、三重态的构造见上一节。

\exercisequestion{QM-C15-S02-E01}{教材正文示范：两自旋 1/2 的换基与逆变换}
\textbf{来源与知识点：}印刷 pp.~413--414，无独立 Example 题号；\relatedtopic{QM-C15-S02-T02}{CG 系数与正交换基}。

\textbf{问题。}给定乘积基顺序 $(\ket{++},\ket{+-},\ket{-+},\ket{--})$，以及耦合基顺序 $(\ket{1,1},\ket{1,0},\ket{1,-1},\ket{0,0})$，写出 CG 换基矩阵及逆变换，并用耦合基表示 $\ket{+-}$。

\begin{checkedsolution}
定义矩阵 $C$ 的各列为耦合基向量在乘积基中的分量。由上一节的四个态得到
\begin{equation}
  C=\begin{pmatrix}
    1&0&0&0\\
    0&1/\sqrt2&0&1/\sqrt2\\
    0&1/\sqrt2&0&-1/\sqrt2\\
    0&0&1&0
  \end{pmatrix},\qquad C^{\mathsf T}C=I.
\end{equation}
若同一个物理态的两套坐标列分别为 $v_{\rm p},v_{\rm c}$，则
\begin{equation}
  \boxed{v_{\rm p}=Cv_{\rm c},\qquad v_{\rm c}=C^{\mathsf T}v_{\rm p}.}
\end{equation}
特别地，$\ket{+-}$ 的乘积基坐标为 $(0,1,0,0)^{\mathsf T}$，所以
\begin{equation}
  \boxed{\ket{+-}=\frac{\ket{1,0}+\ket{0,0}}{\sqrt2}.}
\end{equation}
这里 $C$ 变换的是同一个态的坐标；指定行列顺序后，逆矩阵就是转置。
\end{checkedsolution}

\exercisequestion{QM-C15-S02-E02}{教材习题 15.2.2 第 (1) 问：\texorpdfstring{$\tfrac12\otimes1$}{1/2 与 1 合成}的 CG 系数}
\textbf{来源：}印刷 p.~413 / PDF p.~419。

\textbf{对应知识点：}\relatedtopic{QM-C15-S02-T01}{合成规则}与\relatedtopic{QM-C15-S02-T02}{最高态构造}。

\textbf{题目。}求 $\tfrac12\otimes1=\tfrac32\oplus\tfrac12$ 的 CG 系数。固定第一因子为 $j_1=1/2$、第二因子为 $j_2=1$，记
\begin{equation}
  \ket{\pm,m_2}=\ket{\tfrac12,\pm\tfrac12}\otimes\ket{1,m_2}.
\end{equation}
从 $\ket{\tfrac32,\tfrac32}=\ket{+,1}$ 下降，求 $\ket{\tfrac32,\tfrac12}$；再由正交性、归一化求 $\ket{\tfrac12,\tfrac12}$，取其实系数中 $\ket{+,0}$ 项为正。补齐其余态。

\begin{checkedsolution}
\textbf{1. 下降求第一组。}两套基分别给出
\begin{align}
  J_-\ket{\tfrac32,\tfrac32}&=\hbar\sqrt3\ket{\tfrac32,\tfrac12},\\
  (J_{1-}+J_{2-})\ket{+,1}&=\hbar\ket{-,1}+\hbar\sqrt2\ket{+,0}.
\end{align}
比较得到
\begin{equation}
  \boxed{\ket{\tfrac32,\tfrac12}
     =\sqrt{\frac23}\ket{+,0}+\sqrt{\frac13}\ket{-,1}.}
\end{equation}

\textbf{2. 正交求新的最高态。}设 $\ket{\tfrac12,\tfrac12}=c\ket{+,0}+d\ket{-,1}$。条件为
\begin{equation}
  \sqrt{\frac23}c+\sqrt{\frac13}d=0,\qquad c^2+d^2=1,\qquad c>0.
\end{equation}
所以 $c=\sqrt{1/3}$、$d=-\sqrt{2/3}$，即
\begin{equation}
  \boxed{\ket{\tfrac12,\tfrac12}
     =\sqrt{\frac13}\ket{+,0}-\sqrt{\frac23}\ket{-,1}.}
\end{equation}
它被 $J_+$ 消去，且 $m=1/2$，因此确为 $j=1/2$ 的最高态。

\textbf{3. 补齐全部系数。}继续下降，或用负 $m$ 对称关系，得到下表。$j=3/2$ 的反转因子为 $+1$，$j=1/2$ 的为 $-1$。
\begin{center}
\renewcommand{\arraystretch}{1.4}
\begin{tabular}{@{}ll@{}}
  \toprule
  耦合态 & 乘积基展开 \\
  \midrule
  $\ket{\tfrac32,\tfrac32}$ & $\ket{+,1}$ \\
  $\ket{\tfrac32,\tfrac12}$ & $\sqrt{2/3}\ket{+,0}+\sqrt{1/3}\ket{-,1}$ \\
  $\ket{\tfrac32,-\tfrac12}$ & $\sqrt{1/3}\ket{+,-1}+\sqrt{2/3}\ket{-,0}$ \\
  $\ket{\tfrac32,-\tfrac32}$ & $\ket{-,-1}$ \\
  $\ket{\tfrac12,\tfrac12}$ & $\sqrt{1/3}\ket{+,0}-\sqrt{2/3}\ket{-,1}$ \\
  $\ket{\tfrac12,-\tfrac12}$ & $\sqrt{2/3}\ket{+,-1}-\sqrt{1/3}\ket{-,0}$ \\
  \bottomrule
\end{tabular}
\end{center}
六个态正交归一，$4+2=2\times3$ 完整覆盖原空间。表中的张量积顺序始终为 $1/2\otimes1$；比较其他表格时须同时核对因子顺序与最高态相位。
\end{checkedsolution}

\exercisequestion{QM-C15-S02-E03}{教材习题 15.2.5：三重态与单重态投影算符}
\textbf{来源与知识点：}印刷 p.~415 / PDF p.~421；\relatedtopic{QM-C15-S02-T02}{固定总角动量子空间}。

\textbf{题目。}两个粒子均有自旋 $1/2$，在其联合自旋空间定义
\begin{equation}
  \mathbb P_1=\frac34I+\frac{\mathbf S_1\cdot\mathbf S_2}{\hbar^2},\qquad
  \mathbb P_0=\frac14I-\frac{\mathbf S_1\cdot\mathbf S_2}{\hbar^2}.
\end{equation}
证明 $\mathbb P_i\mathbb P_j=\delta_{ij}\mathbb P_j$，并说明它们分别投影到 $\tfrac12\otimes\tfrac12=1\oplus0$ 中的三重态、单重态子空间。

\begin{checkedsolution}
由 $\mathbf S_1\cdot\mathbf S_2=(S^2-S_1^2-S_2^2)/2$，以及 $S_1^2=S_2^2=3\hbar^2I/4$，
\begin{equation}
  \mathbf S_1\cdot\mathbf S_2\ket{s,m}
    =\frac{\hbar^2}{2}\left[s(s+1)-\frac32\right]\ket{s,m}.
\end{equation}
代入两个可能的 $s$：
\begin{center}
\begin{tabular}{@{}llll@{}}
  \toprule
  子空间 & $\mathbf S_1\cdot\mathbf S_2$ 本征值 & $\mathbb P_1$ 本征值 & $\mathbb P_0$ 本征值 \\
  \midrule
  三重态 $s=1$ & $\hbar^2/4$ & $1$ & $0$ \\
  单重态 $s=0$ & $-3\hbar^2/4$ & $0$ & $1$ \\
  \bottomrule
\end{tabular}
\end{center}
因此 $\mathbb P_1$ 保留三重态、消去单重态，$\mathbb P_0$ 恰好相反。四个耦合态构成完整基，所以在整个空间内
\begin{equation}
  \boxed{\mathbb P_1^2=\mathbb P_1,\quad
    \mathbb P_0^2=\mathbb P_0,\quad
    \mathbb P_1\mathbb P_0=\mathbb P_0\mathbb P_1=0.}
\end{equation}
此外，不同因子的自旋算符对易，故 $\mathbf S_1\cdot\mathbf S_2$ 厄米，进而 $\mathbb P_i^\dagger=\mathbb P_i$。由定义直接得到 $\mathbb P_1+\mathbb P_0=I$。其谱表达为
\begin{equation}
  \boxed{\mathbb P_1=\sum_{m=-1}^1\ketbra{1,m}{1,m},\qquad
    \mathbb P_0=\ketbra{0,0}{0,0}.}
\end{equation}
于是任意态唯一分解为 $\ket\psi=\mathbb P_1\ket\psi+\mathbb P_0\ket\psi$。对归一化态，测得总自旋 $s=i$ 的概率为 $\expect{\mathbb P_i}=\norm{\mathbb P_i\ket\psi}^2$。
\end{checkedsolution}
